\documentclass[a4paper,11pt]{article}
\usepackage[utf8]{inputenc}
\usepackage[T1]{fontenc}
\usepackage[french]{babel}
\usepackage[left=2cm,right=2cm,top=2cm,bottom=2cm]{geometry}
\usepackage{xcolor}
\usepackage{hyperref}
\usepackage{fontawesome5}
\usepackage{parskip}

\definecolor{primary}{RGB}{0, 70, 173}
\hypersetup{colorlinks=true, urlcolor=primary}

\begin{document}

\noindent
\begin{minipage}[t]{0.5\textwidth}
    \textbf{Loïc Aron MBASSI EWOLO}\\
    Élève-Ingénieur Bac+5 -- Double diplôme ENSEA\\[3pt]
    \faMapMarker\ Cergy, France\\
    \faPhone\ (+33) 7 59 52 33 98\\
    \faEnvelope\ \href{mailto:loicaron.mbassiewolo@ensea.fr}{loicaron.mbassiewolo@ensea.fr}
\end{minipage}
\hfill
\begin{minipage}[t]{0.4\textwidth}
    \raggedleft
    Cergy, le 10 mars 2026
\end{minipage}

\vspace{1.2cm}

\hfill
\begin{minipage}[t]{0.5\textwidth}
    \raggedleft
    \textbf{Capgemini}\\
    Service Recrutement\\
    Issy-les-Moulineaux, France
\end{minipage}

\vspace{1cm}

\noindent\textbf{Objet :} Candidature en alternance -- Ingénieur Systèmes

\vspace{0.8cm}

Madame, Monsieur,

Modéliser un drone en conditions nominales et dégradées, spécifier les exigences de sécurité et de performance, puis tracer chaque exigence vers une solution technique (détection de défauts, reconfiguration de contrôleur, gain scheduling) : ce projet à l'ENSEA m'a initié à la démarche d'ingénierie système, de l'analyse des besoins à la validation par simulation. C'est cette approche structurée que je souhaite approfondir au sein de l'équipe Systèmes de Capgemini, dans un contexte automobile et embarqué.

Le challenge CoHoMa (robotique autonome, ENSEA/Armée française) m'a confronté à l'intégration système à l'échelle d'un robot complet : définition de l'architecture matérielle et logicielle, intégration de capteurs hétérogènes (Lidar 3D, caméra de profondeur), spécification des interactions entre composants et containerisation Docker pour l'intégration continue. Mon projet d'architecture microservices GoFind (API Gateway, RabbitMQ, Docker) a renforcé ma capacité à cartographier les flux de données et les dépendances entre sous-systèmes, compétences directement applicables à la structuration d'architectures système et à l'élaboration de SCDR.

En stage à INRIA Saclay (équipe TRIBE, avril à août 2026), je développe un pipeline Python d'analyse automatisée incluant de la traçabilité entre données d'entrée et résultats de scoring, avec rédaction scientifique en anglais. Mon passage au MILA (Institut Québécois d'IA) en environnement anglophone garantit mon aisance à un niveau B2+ dans les contextes internationaux que propose Capgemini.

Je suis disponible dès septembre 2026 pour une alternance de 12 mois. Analyser les besoins, spécifier les fonctions et garantir la traçabilité entre exigences métiers et solutions techniques dans le secteur automobile est le type de mission dans lequel je veux m'investir.

Je reste à votre disposition pour un entretien afin de discuter de ma candidature.

Je vous prie d'agréer, Madame, Monsieur, l'expression de mes salutations distinguées.

\vspace{0.8cm}

\textbf{Loïc Aron MBASSI EWOLO}\\
\textit{Élève-Ingénieur ENSEA / Polytechnique Yaoundé}

\end{document}
